\section{Related Work}
\label{sec:related}

\subsection{Workflow-aware KV cache management}

The closest body of work manages the KV cache of multi-agent workflows using
knowledge of what will run next, and differs mainly in where that knowledge comes
from.

\textbf{KVFlow}~\cite{pan2025kvflow} abstracts the agent execution schedule as an
Agent Step Graph and assigns each agent a steps-to-execution value estimating its
temporal proximity to future activation; that value drives a node-level eviction
policy and a fully overlapped prefetch that loads tensors from CPU to GPU in
background threads for agents scheduled in the next step. Its target is explicitly
the KV corresponding to agents' \emph{fixed} prompts. It reports up to
1.83$\times$ speedup for single workflows with large prompts and up to
2.19$\times$ with many concurrent workflows against SGLang with a hierarchical
radix cache.

\textbf{PBKV}~\cite{zheng2026pbkv} targets \emph{dynamic} workflows, where the
sequence of invoked agents depends on the task. It predicts agent invocations
several steps ahead by fusing guidance from historical workflows with the context
of the target workflow, scores cache entries by reuse potential, and applies the
predictions conservatively in both eviction and prefetching. It reports up to
1.85$\times$ over LRU on dynamic workflows and up to 1.26$\times$ over KVFlow on
the static workflow in its evaluation.

\textbf{CacheScout}~\cite{zhang2026cachescout} removes the offline dependency
entirely: it learns agent execution transitions online, requiring no predefined
workflow graph, no annotations, and no trace replay. It reports that predicting
the next agent as the most frequent observed successor reaches 76--86\% top-1
accuracy within 50 observed dispatches, and that recurring anchor content accounts
for 53--62\% of prompt tokens across its four evaluated workloads. It materializes
predicted context by issuing a lightweight synthetic warmup request containing
only the predicted agent's reusable anchor.

\textbf{Pythia}~\cite{yu2026pythia} goes furthest along the content axis. It
mines logged execution traces into probabilistic automata, identifies each
request through three annotations attached at submission
(\texttt{workflow\_type\_id}, \texttt{workflow\_id}, \texttt{agent\_id}), and
prefills not only shared prefixes but argument content: upcoming prompts are
assembled by injecting inputs and outputs recorded in the execution history,
addressed through strict memory pointers into prior requests, and precomputed
in idle GPU cycles. Its knowledge is the trace: a workflow or a path outside
the history has nothing to instantiate.

\textbf{SAGA}~\cite{saga2026} makes the workflow, rather than the individual
call, the schedulable unit on a GPU cluster. Its agent execution graphs arrive
from framework hints (orchestrator callbacks and message logs) or are inferred
from request traces at 87\% accuracy, with a fallback to request-level
scheduling while a new agent type accumulates history; eviction is ranked by
the predicted reuse of an agent's successors, reported within 1.31$\times$ of
B\'el\'ady's offline optimal on SWE-bench.

\textbf{Helium}~\cite{helium2026} is the nearest neighbor on the compile-time
axis. It treats an agentic workflow as a query plan: the developer composes the
workflow in a dedicated DSL, LLM calls become plan operators, compilation
identifies the static token sequences of each prompt, and the KV for those
sequences is precomputed at the workflow's first execution and reused
afterward, with the plan's leaves partitioned across workers to maximize reuse.
It reports up to 1.56$\times$ over KVFlow. The compiled knowledge stops at
static prefixes: the plan is authored by the developer rather than recovered
from a program, argument values are not speculated, and cache eviction remains
the engine's native LRU.

\subsection{Static analysis of agent programs}

AgentFlow~\cite{wang2026agentflow} recovers agent dependencies from agent programs
into a framework-agnostic Agent Dependency Graph, capturing component, control-flow
and data-flow dependencies, and uses it for governance and security analyses such
as Agent Bill of Materials generation and prompt-to-tool risk detection; it is
implemented for five frameworks and evaluated on a corpus of 5{,}399 real-world
agent programs. It establishes that agent program structure is statically
recoverable and that doing so is worth engineering, but it consumes that structure
for security rather than for serving. \sysname{} shares the premise and applies it
to the runtime's KV schedule.

\subsection{Serving substrates}

\sysname{} builds on established serving machinery rather than replacing it.
vLLM~\cite{kwon2023pagedattention} supplies paged KV management and the automatic
prefix caching used as our baseline; SGLang~\cite{zheng2023sglang} established
radix-structured prefix reuse for structured LM programs;
Sarathi-Serve~\cite{agrawal2024sarathi} contributes the chunked-prefill scheduling
that lets us admit speculative prefill incrementally; and
XGrammar~\cite{dong2025xgrammar} provides the constrained-decoding machinery whose
grammar compilation we hoist into the idle window. The programming model is byLLM,
the runtime of the meaning-typed abstraction introduced by
MTP~\cite{dantanarayana2025mtp}, which is what makes the prompt a compiler artifact
in the first place.
